\documentclass[a4paper,12pt]{article}
\usepackage[utf8]{inputenc}
\usepackage[T1]{fontenc}
\usepackage[spanish]{babel}
\usepackage{amssymb}
\usepackage{geometry}
\geometry{margin=2.2cm}

\title{Checklist de la sesion 2}
\author{Entorno}
\date{}

\begin{document}

\maketitle

\section*{Crear un documento municipal estructurado}

Marca cada punto cuando lo hayas comprobado.

\begin{itemize}
    \item[$\square$] He creado la carpeta \texttt{e2}.
    \item[$\square$] He creado \texttt{informe\_servicios.tex}.
    \item[$\square$] He compilado el documento con el icono de Play.
    \item[$\square$] Se ha generado \texttt{informe\_servicios.pdf}.
    \item[$\square$] He revisado el titulo del PDF.
    \item[$\square$] He comprobado que la lista se lee bien.
    \item[$\square$] He comprobado que la tabla cabe en la pagina.
    \item[$\square$] He hecho una modificacion en el archivo \texttt{.tex}.
    \item[$\square$] He recompilado y visto el cambio en el PDF.
    \item[$\square$] He pedido a Copilot una explicacion sin aceptar cambios automaticamente.
    \item[$\square$] He comprobado que no hay datos personales reales.
\end{itemize}

\section*{Nombre}

\vspace{1cm}

\hrule

\end{document}
